\documentclass[11pt,a4paper]{ctexart}

\usepackage{amsmath,amssymb,bm}
\usepackage{geometry}
\usepackage{enumitem}
\usepackage{fancyhdr}
\usepackage{lastpage}
\usepackage{xcolor}
\usepackage{tcolorbox}
\usepackage{tikz}
\usepackage{titlesec}
\usepackage{microtype}
\usetikzlibrary{arrows.meta,calc}
\usepackage{graphicx} 

\geometry{left=2.15cm,right=2.15cm,top=2.0cm,bottom=2.0cm}
\setlength{\parindent}{2em}
\setlength{\parskip}{0.25em}
\linespread{1.12}

\definecolor{mainblue}{RGB}{39,91,150}
\definecolor{lightblue}{RGB}{232,241,252}
\definecolor{darkgray}{RGB}{65,65,65}
\definecolor{lightgray}{RGB}{245,246,248}

\titleformat{\section}{\large\bfseries\color{mainblue}}{}{0em}{}[\titlerule]
\titleformat{\subsection}{\normalsize\bfseries\color{darkgray}}{}{0em}{}
\titlespacing*{\section}{0pt}{0.9em}{0.45em}
\titlespacing*{\subsection}{0pt}{0.65em}{0.25em}

\setlist[enumerate]{leftmargin=2.2em,itemsep=0.25em,topsep=0.25em}
\setlist[itemize]{leftmargin=2.0em,itemsep=0.2em,topsep=0.2em}

\pagestyle{fancy}
\fancyhf{}
\lhead{\small 格物助航｜每日一练}
\rhead{\small 2026.6.13}
\cfoot{\small 第 \thepage\ 页 / 共 \pageref{LastPage} 页}
\renewcommand{\headrulewidth}{0.4pt}

\newtcolorbox{problemBox}{
  colback=lightblue,
  colframe=mainblue,
  arc=2mm,
  boxrule=0.7pt,
  left=1.0em,
  right=1.0em,
  top=0.8em,
  bottom=0.8em,
  before skip=0.6em,
  after skip=0.8em
}

\newtcolorbox{tipBox}{
  colback=lightgray,
  colframe=darkgray,
  arc=1.5mm,
  boxrule=0.5pt,
  left=0.9em,
  right=0.9em,
  top=0.6em,
  bottom=0.6em,
  before skip=0.5em,
  after skip=0.5em
}

\newcommand{\dd}{\mathrm{d}}
\newcommand{\ii}{\bm{i}}
\newcommand{\jj}{\bm{j}}

\begin{document}

\begin{center}
  {\Large\bfseries\color{mainblue} 格物助航｜每日一练}\par
  \vspace{0.25em}
  {\large 力学第八章：狭义相对论}\par
  \vspace{0.35em}
  {\small 2026.6.13}\par
  \vspace{0.35em}
  {\small 编写：2542 桃酥OvO}
\end{center}

\vspace{-0.3em}

\begin{problemBox}
\noindent
\begin{minipage}[c]{0.50\textwidth}
\textbf{题目}\quad
惯性系$S',S$间的关系如图所示。静长为$l$的直尺$AB$静止于$x'y'$平面，且与$x$轴夹角为$45^\circ$，质点$P$沿杆运动，相对于杆的速率为常量$u$，试在$S',S$系分别计算质点$P$从杆的$A$端到$B$端所经历的时间$T',T$。
\end{minipage}
\hfill
\begin{minipage}{0.45\textwidth}
\centering
\includegraphics[width=0.75\linewidth]{\fig_6.13\graph.jpg}
\end{minipage}
\end{problemBox}


\section*{详细解法}

\subsection*{解法（1）：洛伦兹变换}
显然$T'=\dfrac{l}{u}$；

下面给出三种方法求解$T$，旨在体现相对论体系的自洽性，以及便于从多个角度看待同一问题，便于对比分析。

质点$P$在杆$A$端、$B$端对应的空时坐标依次为$\{x_A',t_A'\},\{x_A,t_A\}$与$\{x_B',t_B'\},\{x_B,t_B\}$。
杆在$S'$系中的坐标差与时间差满足：
\[
x_B'-x_A' = l\cos45^\circ = \frac{\sqrt{2}}{2}l,\quad t_B'-t_A' = T' = \frac{l}{u}.
\]
根据洛伦兹变换可得：
\[
T = t_B - t_A
= \frac{(t_B' - t_A') + \dfrac{v}{c^2}(x_B' - x_A')}{\sqrt{1-\beta^2}}
= \frac{1 + \dfrac{\sqrt{2}}{2}\dfrac{uv}{c^2}}{\sqrt{1-\beta^2}}\cdot\frac{l}{u},\quad \beta = \frac{v}{c}.
\]

\subsection*{解法（2）：$x$方向速度变换}
质点$P$在$x'$、$x$方向的速度分量为：
\[
u_x' = \frac{\sqrt{2}}{2}u,\quad
u_x = \frac{u_x' + v}{1 + \dfrac{u_x'v}{c^2}}
= \frac{\dfrac{\sqrt{2}}{2}u + v}{1 + \dfrac{\sqrt{2}}{2}\dfrac{uv}{c^2}}.
\]

$S$系中杆沿$x$方向发生长度收缩，其$x$方向长度为$\sqrt{1-\beta^2}\,l\cos45^\circ$；$P$相对杆$B$端的$x$方向追击速度为$u_x-v$，因此：
\[
T = \frac{\sqrt{1-\beta^2}\,l\cos45^\circ}{u_x - v}
= \frac{1 + \dfrac{\sqrt{2}}{2}\dfrac{uv}{c^2}}{\sqrt{1-\beta^2}}\cdot\frac{l}{u}.
\]

\subsection*{解法（3）：$y$方向速度变换}
本方法思路与解法（2）相似。质点$P$在$y'$方向速度分量：
\[
u_y' = \frac{\sqrt{2}}{2}u,\quad
u_y = \frac{\sqrt{1-\beta^2}\,u_y'}{1 + \dfrac{u_x'v}{c^2}}
= \frac{\dfrac{\sqrt{2}}{2}\sqrt{1-\beta^2}\,u}{1 + \dfrac{\sqrt{2}}{2}\dfrac{uv}{c^2}}.
\]

$y$方向无长度收缩，由运动学关系得：
\[
T = \frac{l\sin45^\circ}{u_y}
= \frac{1 + \dfrac{\sqrt{2}}{2}\dfrac{uv}{c^2}}{\sqrt{1-\beta^2}}\cdot\frac{l}{u}.
\]

\section*{技巧与知识点}

\begin{tipBox}
\begin{enumerate}
  \item \textbf{知识点1：} 狭义相对论中长度收缩仅发生在两惯性系的相对运动方向（$x$方向），垂直运动方向（$y$方向）物体长度保持不变。
  \item \textbf{知识点2：} 本题综合考查洛伦兹时空变换、相对论速度变换两大核心公式，是狭义相对论运动学典型综合题型。
\end{enumerate}
\end{tipBox}

\section*{复习建议}

1. 熟练背诵洛伦兹坐标变换、时间变换以及相对论速度变换的正、逆公式，明确公式的适用参考系。

2. 严格区分固有长度、运动长度、固有时间、运动时间四个基础概念，解题时优先判定研究对象所处参考系。

\end{document}